\documentclass[
 aip,
 jcp,
 amsmath,amssymb,
 preprint
]{revtex4-2}

\usepackage{graphicx}
\usepackage{dcolumn}
\usepackage{bm}
\usepackage{hyperref}
\usepackage[T1]{fontenc}
\usepackage[utf8]{inputenc}

\begin{document}

\title{From Projected Subspaces to Full-Space Implementations:
Representation-Equivalence Audits for Open-Shell ADAPT-VQE}

\author{Lucia Mal\'{\i}\v{c}kov\'a}
\email{lucy@thesafehaven.ai}
\affiliation{Modelos Inteligencia Artificial S.L., Cl. Tajinaste 54, 386 20 San Miguel de Abona, Santa Cruz de Tenerife, Spain}

\author{Petr Klenovsk\'y}
\email{klenovsky@physics.muni.cz}
\affiliation{Department of Condensed Matter Physics, Faculty of Science, Masaryk University, Kotl\'a\v{r}sk\'a 267/2, 611 37 Brno, Czech Republic}
\affiliation{Czech Metrology Institute, Okru\v{z}n\'i 31, 638 00 Brno, Czech Republic}

\date{\today}

\begin{abstract}
A variational state optimized in a symmetry-projected space need not be the state prepared by exponentiating the corresponding unprojected parent generators.
For an orthogonal target-space projector $P$, $Q=I-P$, and an anti-Hermitian generator $A$, exact equivalence for all amplitudes requires target-space invariance, $QAP=0$.
We formulate this condition as a representation-equivalence audit for open-shell ADAPT-VQE.
For a canonical CAS(15e,9o) type-1 Cu(II) Hamiltonian, 274 of 323 bare Hartree--Fock-origin single, double, and triple excitations violate this condition.
Projected-doublet ADAPT reaches 1.5719 m$E_h$ with 121 operators, whereas the same amplitudes in the corresponding full-space sequence give 464.730 m$E_h$ error, 0.07343 fidelity with the projected state, and 54.45\% quartet weight.
Reoptimizing the same 121 parent generators in the same order under the hard constraint $p_D\ge0.999$ reduces the best converged error to 8.482 m$E_h$, showing that the catastrophic same-angle failure is primarily a representation/parameter-transfer failure but is not completely repaired by full-space reoptimization of the fixed sequence.
An independent 12-qubit NO benchmark exhibits the same mechanism more mildly.
A spin-preserving generalized full-space construction removes the representation ambiguity and reaches 1.493 and 1.181 m$E_h$ at 18 and 24 qubits.
For the 18-qubit positive control, selectively refining only four Suzuki-2 factors and re-transpiling with a fixed compiler seed passes the predeclared circuit-equivalence gate at 19,660 controlled-X (CX) gates and depth 23,016.
Circuit synthesis and final-energy measurement therefore constitute separate validation layers.
The contribution is a validation framework, not a new spin-adapted ADAPT algorithm.
\end{abstract}

\maketitle

\section{Introduction}

Variational quantum eigensolvers (VQEs) and ADAPT-VQE are widely studied approaches to quantum electronic-structure simulation \cite{peruzzo2014vqe,mcclean2016theory,mcardle2020review,grimsley2019adapt}.
Qubit pools, symmetry-aware and symmetry-complete pools, spin projection, compact excitation circuits, and pruning have all been developed subsequently \cite{tang2021qubitadapt,shkolnikov2023roadblocks,bertels2022symmetry,tsuchimochi2022spinproj,magoulas2023cnot,vaquero2025pruned}.
For open-shell systems, fixed particle number and $M_S$ do not determine total spin, and symmetry preservation can be essential for state targeting \cite{gard2020symmetry,seki2020symmetry,seki2022projection,burton2023disco,magoulas2026symmetry}.
Likewise, finite product-formula implementations of spin-adapted generators can introduce symmetry error, motivating exact or closed-form factorizations \cite{magoulas2025closedform,jain2026factorization}.

Several neighboring approaches therefore solve important but distinct problems.
Symmetry projection restores or enforces a desired sector in the variational state or objective \cite{seki2020symmetry,tsuchimochi2022spinproj,seki2022projection}; symmetry-preserving state-preparation circuits and symmetry-aware or symmetry-complete pools constrain the accessible full-space manifold \cite{gard2020symmetry,shkolnikov2023roadblocks,burton2023disco,magoulas2026symmetry}; contextual-subspace VQE deliberately defines a reduced-space Hamiltonian and projected operator pool \cite{weaving2023stabilizer}; and recent exact or closed-form factorizations address how a specified spin-adapted full-space unitary is compiled \cite{magoulas2025closedform,jain2026factorization}.
None of these statements by itself establishes that, for an arbitrary parent generator $A$, the reduced-space unitary $\exp(\theta U^\dagger A U)$ is identical to the target-space compression of the full-space parent evolution $\exp(\theta A)$.
The present work isolates and audits precisely this representation-identification step.

Projected adaptive pools are not new.
For example, contextual-subspace VQE optimizes projected operators directly in a reduced space \cite{weaving2023stabilizer}.
Such algorithms are legitimate reduced-space variational methods.
The narrower issue considered here arises only if the corresponding unprojected parent-generator sequence is subsequently treated as the full-space implementation, or assigned a circuit resource count, without proving that it realizes the same variational state.
Operationally, the ambiguity can enter when (i) parent excitations are projected or compressed into a target sector, (ii) ADAPT selection and parameter optimization are performed in that reduced representation, (iii) the selected parent labels and amplitudes are retained, and (iv) the unprojected parents are later exponentiated to construct a fermionic or qubit circuit or to quote resources.
The reduced calculation never explores the complementary $Q$ sector, so a low reduced-space energy cannot diagnose a nonzero $QAP$.
Figure~\ref{fig:validation_ladder} summarizes the resulting validation ladder.
The underlying invariant-subspace algebra is elementary; the practical question is how large the mismatch can become in an adaptive workflow and how to detect it before reduced-space optimization, full-space evolution, circuit synthesis, and measurement are conflated.

The practical significance is therefore deliberately narrower than a criticism of projected or symmetry-adapted methods themselves.
A reduced-space optimization certifies the variational object defined in that reduced representation.
Associating it with a particular parent full-space fermionic sequence, compiled circuit, or resource estimate is an additional statement that requires its own representation-equivalence check.
We do not assume that prior projected-space studies make this identification; rather, we isolate the condition that must be verified whenever such a connection is made.

\begin{figure*}[htbp]
\centering
\includegraphics[width=0.98\textwidth]{fig1_validation_ladder.pdf}
\caption{Validation layers that should not be identified without an explicit check. The first arrow is exact for all amplitudes only when each implemented parent generator preserves the target subspace, $QA_kP=0$. Subsequent arrows add independent circuit-synthesis and measurement approximations.}
\label{fig:validation_ladder}
\end{figure*}

The primary benchmark is an idealized type-1 (T1, ``blue-copper'') Cu(II) first-shell model, chosen because it provides a chemically motivated open-shell Hamiltonian while retaining exact spin-resolved references \cite{solomon2004bluecopper,solomon2014copper}.
A separate NO radical benchmark and Hubbard-model stress tests probe whether the same representation issue persists beyond that Hamiltonian.
No chemically converged prediction for blue-copper proteins, NO, or OH is claimed.

\section{Model and Validation Strategy}

The canonical Cu Hamiltonian is defined in a complete active space, CAS(15e,9o), comprising 15 active electrons in 9 active spatial orbitals and mapped to 18 qubits.
Its STO-3G orbitals were generated from restricted open-shell Hartree--Fock (ROHF), followed by atomic-valence active-space (AVAS) selection targeting the Cu $3d$ and short Cu--S donor $3p$ manifolds and a spin-constrained complete-active-space self-consistent-field (CASSCF) calculation; the exact geometry and numerical settings are given in the supplementary material \cite{sun2018pyscf,sayfutyarova2017avas}.
An archived CAS(21e,12o) FCIDUMP is used only as a fixed 24-qubit algorithmic benchmark because its dedicated 12-orbital selection driver was not preserved.
The independent NO test uses ROHF/STO-3G CAS(7e,6o) at $R_{\rm NO}=1.1508$~\AA{}.

For the target doublet in the fixed-$M_S=1/2$ sector,
\begin{equation}
\mathcal H_D=\ker\!\left(\hat S_+\big|_{\mathcal H_{M_S=1/2}}\right),
\qquad
P=UU^\dagger,\quad Q=I-P,
\end{equation}
where the columns of $U$ form an orthonormal basis of $\mathcal H_D$.
For a full-space anti-Hermitian fermionic generator $A$, let
\begin{equation}
A_D=U^\dagger A U.
\end{equation}
The lifted reduced-space unitary and the compressed full-space evolution satisfy
\begin{equation}
Ue^{\theta A_D}U^\dagger
=
Pe^{\theta A}P
\quad\text{for all real }\theta
\end{equation}
if and only if
\begin{equation}
QAP=0.
\label{eq:invariance}
\end{equation}
For anti-Hermitian $A$, Eq.~\ref{eq:invariance} is equivalent to $[A,P]=0$.
If invariance fails,
\begin{align}
Qe^{\theta A}P&=\theta QAP+\mathcal O(\theta^2),\\
Pe^{\theta A}P-e^{\theta PAP}P
&=
-\frac{\theta^2}{2}(QAP)^\dagger(QAP)+\mathcal O(\theta^3).
\end{align}
The supplementary material gives the proof, finite-amplitude bounds, and a state-specific telescoping certificate.
The exact projector is used only as a classical validation oracle for these small benchmarks. For the $M_S$-preserving generators considered here, $[A,\hat S^2]=0$ is a projector-free sufficient condition for preserving the target total-spin sector.

The corrected ansatz uses full-space anti-Hermitian generators $B_k$ whose target-space leakage and normalized $S^2$ commutator defect vanish to numerical precision,
\begin{equation}
|\Psi(\boldsymbol\theta)\rangle
=
e^{\theta_M B_M}\cdots e^{\theta_1B_1}|\Phi_{\rm HF}\rangle.
\end{equation}
Pool construction, ADAPT selection, repeated operator use, optimization, pruning, circuit synthesis, and measurement allocation are documented in the SI.
Throughout, 1.6 m$E_h$ is only an internal benchmark for a specified active-space Hamiltonian, not a claim of chemical accuracy for the underlying molecular model.

\section{Results and Discussion}

\subsection{Cu: representation failure and fixed-sequence reoptimization}

The regenerated 121-operator projected-doublet ADAPT sequence reaches
$E=-2518.987495464~E_h$, 1.5719 m$E_h$ above the exact target doublet.
Using the same selected parent generators and the same amplitudes in the full fixed-$M_S$ space instead gives 464.730 m$E_h$ error, $\langle S^2\rangle=2.383537$, doublet weight 0.45549, and only 0.07343 fidelity with the projected state.
Post-projecting that full-space state does not repair the result: its normalized doublet component remains 465.164 m$E_h$ above the target and has fidelity 0.16121 with the projected state.

This mismatch is not caused by one exceptional operator.
Of 323 bare Hartree--Fock-origin single, double, and triple excitations, 274 violate Eq.~\ref{eq:invariance}; the selected sequence contains 114 non-invariant factors.
The supplementary material sequence certificates are correspondingly non-informative for this severe case.

To distinguish a failure of parameter transfer from a failure of the fixed parent sequence itself, we next reoptimized all 121 amplitudes in the \emph{same generator order} directly in the full fixed-$M_S$ space while enforcing
\begin{equation}
p_D=\|U^\dagger\Psi\|^2\ge0.999.
\label{eq:hardspin}
\end{equation}
Starting from zero amplitudes and three independent Gaussian perturbations of that start, all constrained Sequential Least Squares Programming (SLSQP) runs converged; the best three solutions are indistinguishable at 8.4815 m$E_h$, and the four-run range is 8.4815--8.9734 m$E_h$.
The best state has $p_D=0.999$, $\langle S^2\rangle=0.753$, and fidelity 0.97989 with the original projected state.
Its normalized doublet component is still 7.335 m$E_h$ above the exact doublet.
Tightening the constraint to $p_D\ge0.9999$ gives 11.270 m$E_h$ from the zero-amplitude start.
Separate starts near the original projected amplitudes converge to poorer local minima, demonstrating nonconvexity; no global-optimality claim is made.

We also optimized the same fixed sequence directly for overlap with the projected state.
Starting from the best energy-constrained basin, unconstrained fidelity maximization converges reproducibly to
$F_{\rm proj}=0.988268$ with $p_D=0.99716$.
Imposing $p_D\ge0.999$ during the fidelity maximization gives
$F_{\rm proj}=0.987416$ from three tested starts.
These are converged local optima rather than certified global maxima, but none of the tested full-space searches reproduces the projected state exactly.
The principal same-sequence comparisons are summarized in Table~\ref{tab:parent_reopt}.

\begin{table}[htbp]
\centering
\small
\caption{Canonical Cu representation test for the same 121 parent generators. The constrained rows reoptimize all amplitudes in the full fixed-$M_S$ space without changing operator order.}
\label{tab:parent_reopt}
\begin{tabular}{lrrrr}
\toprule
Construction & $\Delta E$ (m$E_h$) & $p_D$ & $F_{\rm proj}$ & Post-proj.\ $\Delta E_D$ (m$E_h$)\\
\midrule
Projected reduced-space ansatz & 1.5719 & 1 & 1 & 1.5719\\
Same amplitudes, full space & 464.730 & 0.45549 & 0.07343 & 465.164\\
Full-space energy reopt., $p_D\ge0.999$ & 8.4815 & 0.99900 & 0.97989 & 7.3355\\
Full-space fidelity reopt., $p_D\ge0.999$ & 10.8166 & 0.99900 & 0.98742 & 9.8863\\
Full-space energy reopt., $p_D\ge0.9999$ & 11.2696 & 0.99990 & 0.97017 & 11.0305\\
\bottomrule
\end{tabular}
\end{table}

Thus the catastrophic 464.7 m$E_h$ failure is primarily a representation/parameter-transfer failure: the same parent sequence can recover much of the target state after full-space spin-constrained reoptimization.
However, in the tested order it still does not reproduce the 1.6 m$E_h$ projected result.
This distinction is central: the audit does not show that projected-pool methods are invalid, nor that the parent sequence is globally incapable of representing the target; it shows that the projected optimization and the unprojected same-parameter implementation are different variational constructions.

Simple invariant filtering is not an adequate repair.
Only 49 of the 323 Cu excitations are individually invariant.
Using all 49 once gives 210.619 m$E_h$ error, and allowing repeated selection through 100 applications gives 167.340 m$E_h$.
At the Hartree--Fock reference, the projected bare S/D/T pool spans all 239 real tangent directions of the 240-dimensional doublet space, whereas the invariant-only subset spans 49 and the generalized spin-preserving full-space rank-1/rank-2 pool spans 99.
These are local first-order diagnostics, not controllability statements.
The divergence of the projected, full-space, and post-projected trajectories over all 121 prefixes is shown in Fig.~\ref{fig:prefix}.

\begin{figure}[htbp]
\centering
\includegraphics[width=0.92\linewidth]{fig2_projected_prefix_divergence.pdf}
\caption{Canonical Cu negative control. The projected energy approaches the target while the same amplitudes in the parent full-space sequence do not. Post-projection of the resulting full-space state does not recover the optimized projected trajectory.}
\label{fig:prefix}
\end{figure}

\subsection{Independent NO test}

For NO CAS(7e,6o), the exact target doublet lies 250.051 m$E_h$ below the lowest quartet, so representation failure cannot be attributed to collapse into a lower wrong-spin state.
Nevertheless, 225 of 231 bare S/D/T excitations are non-invariant.
Projected ADAPT reaches 1.4286 m$E_h$ with 16 operators; the same parent sequence and amplitudes give 24.537 m$E_h$, doublet weight 0.97211, and fidelity 0.97205 with the projected state.
Here the state-specific SI certificate is nontrivial, guaranteeing $F\ge0.8473$.
Unlike Cu, post-projection almost repairs the result: the post-projected error is 1.4780 m$E_h$ and its fidelity with the projected state is 0.999943.
The same mechanism therefore occurs in an independent molecular Hamiltonian, but its severity is strongly system and sequence dependent.

Only six NO parent excitations are individually invariant and this filtered pool remains 52.154 m$E_h$ above the target.
A generalized spin-preserving pool instead gives a 19-application state at 1.5790 m$E_h$, unit doublet weight, and 0.998967 overlap with the exact lowest-doublet manifold.
Additional OH and Hubbard stress tests are reported in the supplementary material.

\subsection{Representation-faithful full-space ans\"atze}

The generalized spin-preserving pool used below is a representation-faithful \emph{positive control}, not a claim of a new or preferred scalable spin-adapted ADAPT algorithm.
Existing work already provides symmetry-preserving state-preparation circuits, spin-adapted pools, symmetry-completeness analyses, and exact spin-adapted factorizations \cite{gard2020symmetry,shkolnikov2023roadblocks,burton2023disco,magoulas2026symmetry,magoulas2025closedform,jain2026factorization}.
Its purpose here is narrower: to demonstrate on the same benchmarks that, once every implemented generator preserves the target sector, the full-space and reduced-space evolutions do represent the same variational construction to numerical precision.

For Cu, the generalized full-space spin-preserving pools contain 1,080 generators at 18 qubits and 3,762 at 24 qubits.
After repeated selection, reoptimization, and pruning, the stored 18q and 24q ans\"atze contain 50 and 61 applications and reach 1.493 and 1.181 m$E_h$, respectively.
Independent full-space and reduced-space reconstructions agree to machine precision over every prefix.
The final representation-faithful Cu ans\"atze are summarized in Table~\ref{tab:main}.

\begin{table}[htbp]
\centering
\small
\caption{Compact representation-faithful Cu ans\"atze. The 24q case is a fixed-Hamiltonian algorithmic benchmark, not a chemical convergence study.}
\label{tab:main}
\begin{tabular}{lcc}
\toprule
& 18 qubits & 24 qubits\\
\midrule
Active space & CAS(15e,9o) & CAS(21e,12o)\\
Operator applications & 50 & 61\\
Unique generators & 39 & 41\\
$\Delta E$ (m$E_h$) & 1.4927 & 1.1813\\
Exact-doublet fidelity & 0.9826 & 0.9907\\
Residual ($E_h$) & 0.0480 & 0.0585\\
$\langle S^2\rangle$ & 0.7500 & 0.7500\\
\bottomrule
\end{tabular}
\end{table}

The spin-preserving pool basis is not mathematically unique inside multidimensional null spaces.
A normalization-controlled supplementary-material audit finds modest top-group sensitivity along the stored 18q trajectory: random normalized basis changes affect the highest-gradient group at only 4 of 51 sampled prefixes, while a basis-invariant group-gradient score agrees with the canonical choice at 49--50 of 51 prefixes depending on normalization metric.
No basis-independent minimality claim is made for the compact sequences.

\subsection{Circuit and measurement validation}

Second-order Suzuki synthesis introduces a separate approximation.
The 0.1 m$E_h$ synthesis-error threshold was fixed before the selective-refinement search so that circuit-synthesis error would contribute at most 6.25\% of the 1.6 m$E_h$ model-space accuracy benchmark, leaving the ansatz error as the dominant contribution.
For the 18q Cu state, four uniform repetitions give 0.135 m$E_h$ synthesis error and circuit-to-fermionic fidelity 0.999484 at 61,482 CX gates and abstract-gate-set depth 71,636, so this uniform circuit narrowly fails the predeclared $0.1$ m$E_h$ synthesis-error gate.
The factorization error is highly nonuniform, however.
A restricted emulator search identifies the schedule $r_9=4$, $r_{21}=r_{25}=r_{34}=2$, with all other applications at $r=1$.
Re-transpiling exactly this schedule with Qiskit 2.2.3, optimization level 1, the abstract basis $\{R_z,\sqrt X,X,\mathrm{CX}\}$, and fixed \texttt{seed\_transpiler}=9272026 gives
$|\delta E_{\rm synth}|=0.095021$ m$E_h$,
$F_{\rm circuit,fermionic}=0.9995825$,
$\langle S^2\rangle=0.750724$,
and target-$M_S$ weight $0.9999999999994$.
It therefore passes all four predeclared state-level gates at 19,660 CX gates and depth 23,016, reductions of about 68.0\% and 67.9\%, respectively, relative to uniform $r=4$.
For 24q, one repetition gives 0.0906 m$E_h$ and fidelity 0.999694 at 22,574 CX gates and depth 26,409; for NO, one repetition gives 0.00293 m$E_h$ and fidelity 0.9999939.
All reported resources are abstract-gate-set compilations without a device coupling map and are not optimal-circuit or hardware-resource bounds.
Further schedule and provenance details are given in the supplementary material.

Final-energy measurement is likewise protocol dependent.
For the chosen greedy qubit-wise commuting (QWC) groupings and ideal state preparation, variance-optimal allocation requires approximately $1.43\times10^8$ shots (18q) and $5.46\times10^8$ shots (24q) for a 1.6 m$E_h$ statistical uncertainty.
These numbers exclude ADAPT gradient measurements, optimizer evaluations, device noise, and alternative measurement strategies.

\subsection{Scope of claims}

Spin projection, symmetry-preserving state preparation, projected pools, generalized excitations, repeated selection, spin-adapted or symmetry-complete pools, and pruning are established ideas \cite{gard2020symmetry,seki2020symmetry,tsuchimochi2022spinproj,seki2022projection,weaving2023stabilizer,shkolnikov2023roadblocks,burton2023disco,magoulas2026symmetry}.
Our contribution is not a new projector, a new general spin-adapted pool, or a new factorization method, but a representation-equivalence validation framework for connecting a reduced-space optimization to a parent full-space implementation or resource estimate.
The exact projector used here serves only as a small-system validation oracle.

The molecular calculations are fixed algorithmic benchmarks rather than chemically converged predictions.
The compact ADAPT sequences are search-history dependent, and the reported circuit and measurement resources are specific to the documented synthesis, transpilation, and measurement protocols; no global minimality, hardware-optimality, or quantum-advantage claim is made.

\section{Conclusions}

A projected adaptive ansatz and the corresponding parent full-space sequence are guaranteed to define the same generator-by-generator variational construction for arbitrary amplitudes when every implemented parent generator preserves the target space.
The converse necessity applies to this amplitude-independent representation guarantee; accidental agreement at isolated parameter values or cancellations within one particular finite sequence are not excluded.
The canonical Cu example shows how severely this identification can fail: 1.5719 m$E_h$ in projected space becomes 464.730 m$E_h$ when the same amplitudes are applied to the parent sequence.
Full-space energy reoptimization of that fixed sequence under $p_D\ge0.999$ reduces the best converged error to 8.482 m$E_h$ and raises the projected-state fidelity to 0.97989.
Direct fidelity maximization under the same purity constraint reaches 0.98742 but still does not reproduce the projected state.
The catastrophic failure is therefore not evidence that the parent sequence is intrinsically useless; it is primarily evidence that reduced-space parameters cannot be transferred without an equivalence check.
The remaining energy and fidelity gaps show that the fixed-order full-space construction does not completely recover the projected construction in the tested searches, without establishing a global impossibility result.

An independent NO benchmark reproduces the mechanism more mildly, while representation-faithful spin-preserving Cu and NO ans\"atze provide positive controls.
For the 18q Cu positive control, a selectively refined Suzuki-2 schedule also passes the predeclared circuit-equivalence gate after an actual Qiskit re-transpilation, with 19,660 CX gates and depth 23,016 in the abstract basis.
Circuit synthesis and measurement therefore remain separate approximations and resource layers rather than consequences of fermionic-level accuracy.

The practical recommendation is correspondingly narrow:
when a projected or otherwise reduced-space adaptive optimization is connected to a full-space implementation, the optimized state, the parent fermionic sequence, the compiled circuit, and the measurement protocol should be validated as distinct objects.
A low reduced-space energy alone does not establish that a quoted full-space circuit implements the same variational state.

\section*{SUPPLEMENTARY MATERIAL}
See the supplementary material for the complete Hamiltonian-generation protocol; proofs and finite-amplitude bounds; numerical tolerances; state-specific certificates; fixed-parent energy and fidelity reoptimization; Cu prefix and pool audits; tangent-space accessibility; normalization-controlled basis-sensitivity tests; NO, OH, and Hubbard benchmarks; spin-preserving pool construction; circuit and Suzuki diagnostics; measurement grouping; one-particle diagnostics; optimizer robustness; provenance; and excluded calculations.

\begin{acknowledgments}
L.M. gratefully acknowledges the use of the quantum system Euro-Q-Exa, co-funded by the EuroHPC JU, BMFTR, and the Bavarian State Ministry of Science and the Arts, operated by the Leibniz Supercomputing Centre (LRZ) in Garching, Germany, for providing the computational resources for this work. Access to the Euro-Q-Exa infrastructure was awarded under the EuroHPC project proposal No.~EHPC-QCP-2026Q01-017.

P.K. acknowledges support from the European Innovation Council Pathfinder programme under Grant Agreement No.~101185617 (QCEED); the Institutional Subsidy for the Long-Term Conceptual Development of a Research Organization granted to the Czech Metrology Institute by the Ministry of Industry and Trade of the Czech Republic; and the project \emph{Quantum Materials for Applications in Sustainable Technologies}, CZ.02.01.01/00/22\_008/0004572.
\end{acknowledgments}

\section*{AUTHOR DECLARATIONS}

\subsection*{Conflict of Interest}
The authors have no conflicts to disclose.

\subsection*{Author Contributions}
Lucia Mal\'{\i}\v{c}kov\'a: Conceptualization (lead; selection of the research topic); Investigation (lead; execution of the quantum-computing calculations); Writing -- original draft (supporting).

Petr Klenovsk\'y: Formal analysis (lead; many-body calculations outside the quantum-computing workflow); Validation (lead; physical consistency checks); Investigation (supporting); Writing -- original draft (lead); Writing -- review \& editing (lead).

\section*{DATA AVAILABILITY}
The data that support the findings of this study are openly available in the public GitHub repository \url{https://github.com/lucia-malickova/New-VQE-validation}. The repository contains the Hamiltonians, checkpoints, representation-equivalence audits, fixed-parent energy- and fidelity-reoptimization tests, independent benchmarks, circuit outputs, and figure source data used in this work.

\bibliography{references}

\end{document}
